\documentclass[conference]{IEEEtran}
\IEEEoverridecommandlockouts

% Packages
\usepackage{cite}
\usepackage{amsmath,amssymb,amsfonts}
\usepackage{graphicx}
\usepackage{textcomp}
\usepackage{hyperref}
\usepackage{booktabs}

% Figure path
\graphicspath{{../results/paper_figures/}}

\begin{document}

\title{Reinforcement Learning Project Report:\\Implementing and Comparing PPO and DQN on Crafter}

\author{
    \IEEEauthorblockN{Anand Patel}
    \IEEEauthorblockA{\textit{School of Computer Science} \\
    \textit{University of the Witwatersrand} \\
    anand.patel@students.wits.ac.za}
    \and
    \IEEEauthorblockN{[Partner Name]}
    \IEEEauthorblockA{\textit{School of Computer Science} \\
    \textit{University of the Witwatersrand} \\
    [partner.email]@students.wits.ac.za}
}

\maketitle

\begin{abstract}
This report documents our implementation of two reinforcement learning algorithms for the course project. We built PPO (Proximal Policy Optimization) and DQN (Deep Q-Network) agents to play Crafter, a 2D survival game. Starting from a weak baseline, we improved PPO from 5.08\% to 8.61\% through better hyperparameters, adding curiosity-driven exploration, and using a larger network. We tried 11 different approaches, with 7 failing to improve performance. This report explains what we implemented, what worked, what didn't, and what we learned about reinforcement learning in sparse-reward environments.
\end{abstract}

\begin{IEEEkeywords}
Reinforcement Learning, PPO, DQN, Crafter, Project Report
\end{IEEEkeywords}

\section{Introduction}

\subsection{Project Goal}

For this project, we implemented two reinforcement learning algorithms from scratch:
\begin{itemize}
    \item Proximal Policy Optimization (PPO) - Anand's implementation
    \item Deep Q-Network (DQN) - [Partner]'s implementation
\end{itemize}

We trained both agents on Crafter and systematically tested improvements to understand what works for sparse-reward problems.

\subsection{About Crafter}

Crafter is a 2D survival game where the agent must gather resources, craft tools, and survive enemies. It's challenging because:

\begin{itemize}
    \item The agent only sees 64×64 RGB images
    \item Rewards are very sparse (only ~1\% of actions get rewards)
    \item Many tasks need multiple steps (e.g., get wood → build table → craft pickaxe → mine coal)
    \item There are 22 different achievements to unlock
\end{itemize}

The Crafter Score is our main metric—it's the geometric mean of how often the agent unlocks each achievement. Higher scores mean the agent learned more skills.

\subsection{The Two Algorithms}

\textbf{PPO} learns a policy that outputs probabilities for each action. It's stable because it prevents big policy changes.

\textbf{DQN} learns which actions are best by estimating their values. It's sample-efficient because it reuses past experiences.

\subsection{Report Organization}

Section II describes PPO implementation and results. Section III covers DQN. Section IV compares both approaches.

\section{PPO Implementation and Results}

\subsection{How PPO Works}

PPO learns a policy $\pi_\theta(a|s)$ that maps what the agent sees to action probabilities. The key idea is the clipped objective function:

\begin{equation}
L(\theta) = \min(r_t \cdot A_t, \text{clip}(r_t, 0.8, 1.2) \cdot A_t)
\end{equation}

where $r_t$ is how much the new policy differs from the old one, and $A_t$ is the advantage (how good an action was compared to average).

The clipping keeps updates small and prevents the agent from changing too drastically, which makes training stable.

\subsection{Our Implementation}

\textbf{Neural Network Architecture:}

We used an actor-critic network with shared convolutional layers:
\begin{itemize}
    \item 3 convolutional layers to process the 64×64×3 images
    \item Actor head: outputs probabilities for 17 actions
    \item Critic head: outputs a value estimate
    \item Tested 512 and 1024 hidden dimensions
\end{itemize}

\textbf{Training Setup:}
\begin{itemize}
    \item Collected 2048 steps of gameplay
    \item Trained on this data for 10 epochs
    \item Used mini-batches of 64 samples
    \item Ran for 1 million total steps per experiment
    \item Each run took about 2.5 hours on M4 MacBook Pro
\end{itemize}

\subsection{Evaluation 1: Baseline Results}

\textbf{Our Approach:} We started with an initial configuration using common PPO values to establish baseline performance.

\textbf{Settings:}
\begin{itemize}
    \item Learning rate: 0.001
    \item Entropy coefficient: 0.0001
\end{itemize}

\textbf{Results:}
\begin{itemize}
    \item Crafter Score: \textbf{5.08\%}
    \item The agent learned to gather wood (90\%) and saplings (85\%)
    \item It could build tables (75\%)
    \item BUT: Never crafted tools, never found coal, died quickly (165 step episodes)
\end{itemize}

\textbf{Why it struggled:} The learning rate was too high (caused instability) and entropy was too low (agent didn't explore enough).

\subsection{Evaluation 2: Fixing Hyperparameters}

\textbf{What we changed:} Used standard values from published implementations (CleanRL, Stable-Baselines3):
\begin{itemize}
    \item Learning rate: 0.0005 (2× slower)
    \item Entropy coefficient: 0.001 (10× higher)
\end{itemize}

\textbf{Results:}
\begin{itemize}
    \item Crafter Score: \textbf{7.10\%} (+39.8\% improvement)
    \item Wood pickaxe: 51\% (was 40\%)
    \item Wood sword: 48\% (was 35\%)
    \item Started defeating zombies: 48\% (was 30\%)
\end{itemize}

\textbf{Why it worked:} The slower learning rate kept training stable. Higher entropy kept the agent trying new things instead of getting stuck repeating the same actions.

\subsection{Evaluation 3: Adding Curiosity}

\textbf{The Problem:} With sparse rewards, the agent rarely discovers new achievements by accident. It needs help exploring.

\textbf{Our Solution:} We added an Intrinsic Curiosity Module (ICM) that gives the agent bonus rewards for experiencing surprising situations.

\textbf{How ICM Works:}
\begin{enumerate}
    \item A neural network tries to predict what will happen next
    \item When the prediction is wrong (surprising situation), give reward
    \item The agent learns to seek out these surprising situations
    \item Eventually it finds rare achievements this way
\end{enumerate}

The agent now gets two types of rewards:
\begin{itemize}
    \item 80\% from actual achievements (extrinsic)
    \item 20\% from curiosity (intrinsic)
\end{itemize}

\textbf{Results:}
\begin{itemize}
    \item Crafter Score: \textbf{8.27\%} (+16.5\% improvement)
    \item Found coal: 1\% (was 0\% - rare achievement!)
    \item Defeated skeleton: 3\% (was 0\% - very rare!)
    \item Better zombie combat: 54\% (was 48\%)
    \item Survived longer: 200 step episodes (was 182)
\end{itemize}

\textbf{Why it worked:} The curiosity bonus helped systematic exploration. The agent naturally explored early (high curiosity) then focused on achievements later (curiosity decreased).

\begin{figure}[h!]
\centering
\includegraphics[width=0.48\textwidth]{fig6_icm_improvement2.pdf}
\caption{ICM reward analysis for Evaluation 3. Top left: intrinsic (curiosity) vs extrinsic (achievement) rewards over training. Top right: curiosity/achievement ratio showing exploration-to-exploitation transition. Bottom: curiosity decay (agent learns) and achievement growth (agent succeeds). The 0.508 ratio shows good balance.}
\label{fig:icm_eval3}
\end{figure}

\subsection{Evaluation 4: Bigger Network + Less Curiosity}

\textbf{What we tried:} We combined two ideas:
\begin{enumerate}
    \item Made the network bigger (512 → 1024 neurons)
    \item Reduced curiosity slightly (20\% → 15\%)
\end{enumerate}

\textbf{Why we thought this would work:}
\begin{itemize}
    \item Bigger network can remember more complex strategies
    \item Less curiosity means more focus on actual goals
\end{itemize}

\textbf{Results:}
\begin{itemize}
    \item Crafter Score: \textbf{8.61\%} (+4.1\% improvement)
    \item \textbf{Total improvement: 5.08\% → 8.61\% = +69.5\%}
    \item Wood pickaxe: 58\%, Wood sword: 61\%
    \item Coal: 2\% (doubled!)
\end{itemize}

\textbf{Surprising finding:} During training, the average reward was actually lower (5.14) than Evaluation 3 (6.27). But when we evaluated the final agent, it performed better! This taught us that training numbers don't always predict final performance.

\begin{figure}[h!]
\centering
\includegraphics[width=0.48\textwidth]{fig7_icm_improvement3.pdf}
\caption{ICM reward analysis for Evaluation 4 with lower curiosity (β=0.15). The 0.483 ratio is slightly lower than Evaluation 3, showing more focus on achievements. Despite lower training rewards, this configuration achieved the best final score (8.61\%).}
\label{fig:icm_eval4}
\end{figure}

\subsection{Things That Didn't Work}

We tried 7 other strategies that failed. Documenting failures is important for learning:

\subsubsection{RND Curiosity (6.11\%)}

We tried a simpler curiosity method called Random Network Distillation. We attempted this 3 times with different fixes, but all failed.

\textbf{What went wrong:} The curiosity rewards were way too high (42.28 vs ICM's 3.08). The agent got addicted to exploring and never focused on achievements—a "curiosity trap."

\textbf{Lesson learned:} Simpler isn't always better. RND's random features didn't capture what's actually important in Crafter.

\subsubsection{Too Much Curiosity (6.38\%)}

We tried increasing curiosity from 20\% to 30\%.

\textbf{What went wrong:} Another curiosity trap. The agent spent all its time exploring instead of achieving goals.

\textbf{Lesson learned:} Balance is critical. Too much curiosity is as bad as too little.

\subsubsection{Too Much Randomness (4.86\% and 6.08\%)}

We tried higher entropy coefficients (0.005 and 0.002 instead of 0.001).

\textbf{What went wrong:} The agent's policy stayed too random. It never settled on good strategies.

\textbf{Lesson learned:} The agent needs to commit to strategies it learns. Too much randomness prevents this.

\subsubsection{Slower Learning (8.11\%)}

We tried a learning rate of 0.0001 (5× slower).

\textbf{What went wrong:} Got close to baseline (8.11\% vs 8.27\%) but learned too slowly for 1 million steps.

\textbf{Lesson learned:} Learning rate needs to match training time budget.

\subsubsection{Training Longer (7.52\%)}

We tried 1.5 million steps instead of 1 million (50\% more training).

\textbf{What went wrong:} Performance got worse! Dropped from 8.27\% to 7.52\%.

\textbf{Lesson learned:} More training doesn't always help. PPO can degrade because it keeps learning from old experiences that are no longer relevant.

\subsubsection{Dual-Clip PPO (5.79\%)}

We tried a PPO variant that clips differently for positive vs negative advantages.

\textbf{What went wrong:} Made training unstable when combined with curiosity.

\textbf{Lesson learned:} Test modifications carefully. Dual-clip works in some situations but not ours.

\subsubsection{Just Bigger Network (8.36\%)}

We tried only increasing network size without changing curiosity.

\textbf{What happened:} This actually worked (8.36\%) but not as well as combining both changes (8.61\%).

\textbf{Lesson learned:} Combining improvements can work better than individual changes.

\begin{figure}[h!]
\centering
\includegraphics[width=0.48\textwidth]{fig4_failed_attempts.pdf}
\caption{The 7 strategies that didn't improve performance. RND curiosity, high entropy, and extended training all scored worse than our best ICM-based approach. This shows the importance of systematic experimentation.}
\label{fig:failed}
\end{figure}

\subsection{Summary of PPO Results}

\begin{figure}[h!]
\centering
\includegraphics[width=0.48\textwidth]{fig1_score_progression.pdf}
\caption{All 11 PPO experiments we tried. Green dots show successful improvements (5.08\% → 7.10\% → 8.27\% → 8.61\%). Red X's show failed attempts. Numbers show percentage improvements between successful runs.}
\label{fig:all_attempts}
\end{figure}

Figure \ref{fig:all_attempts} shows all 11 attempts. We learned:

\begin{enumerate}
    \item Use standard hyperparameters from the literature as a starting point
    \item Curiosity helps a lot for sparse rewards, but balance is critical
    \item Bigger networks can help, especially for complex tasks
    \item More training isn't always better
    \item Documenting failures helps understand what doesn't work
\end{enumerate}

\section{DQN Implementation and Results}

\textit{[Partner: Fill in this section following the same style as Section II]}

\textit{Use subsections like:}
\begin{itemize}
    \item How DQN Works
    \item Our Implementation
    \item Evaluation 1: Baseline
    \item Evaluation 2: [Your improvement]
    \item Evaluation 3: [Your improvement]
    \item Evaluation 4: [Your improvement]
    \item Things That Didn't Work
    \item Summary of DQN Results
\end{itemize}

\section{Comparing PPO and DQN}

\subsection{Final Scores}

Table \ref{tab:results} shows our final results:

\begin{table}[h]
\centering
\caption{Final Results Comparison}
\label{tab:results}
\begin{tabular}{lcc}
\toprule
\textbf{Metric} & \textbf{PPO} & \textbf{DQN} \\
\midrule
Final Score & 8.61\% & [Partner] \\
Improvement & +69.5\% & [Partner] \\
Training Time & 2.5 hours & [Partner] \\
Wood Pickaxe & 58\% & [Partner] \\
Wood Sword & 61\% & [Partner] \\
Coal (rare) & 2\% & [Partner] \\
\bottomrule
\end{tabular}
\end{table}

\subsection{Exploration Strategies}

The biggest difference between the algorithms was exploration:

\textbf{PPO used curiosity rewards:}
\begin{itemize}
    \item Gave dense rewards for novel experiences
    \item Found rare achievements (coal, skeleton)
    \item But needed careful tuning to avoid curiosity traps
\end{itemize}

\textbf{DQN used [Partner: describe your exploration]:}
\textit{[Partner: Compare your exploration strategy and results]}

\subsection{What We Learned}

\textbf{About sparse rewards:} Standard exploration methods struggle. Curiosity-driven exploration helped PPO find rare achievements that would never be discovered by random chance.

\textbf{About hyperparameters:} Literature values worked much better than our initial guesses. It's worth checking published implementations for starting points.

\textbf{About experimentation:} We tried 11+ different strategies. Many failed, but failures taught us important lessons about what doesn't work.

\textbf{About training:} Training metrics during learning don't always predict final performance. Always evaluate the final model separately.

\subsection{Which Algorithm Worked Better?}

\textit{[Complete after partner finishes DQN section]}

For Crafter specifically:
\begin{itemize}
    \item \textit{[Which got higher score]}
    \item \textit{[Why one performed better]}
    \item \textit{[Trade-offs between approaches]}
\end{itemize}

\subsection{Project Limitations}

\begin{itemize}
    \item Only tested on one environment (Crafter)
    \item Limited compute (~30-40 hours total)
    \item Manual hyperparameter tuning
    \item Could not test more advanced methods (hierarchical RL, etc.)
\end{itemize}

\subsection{If We Had More Time}

\begin{itemize}
    \item Test on other sparse-reward environments
    \item Try combining best parts of both algorithms
    \item Automated hyperparameter search
    \item Test hierarchical approaches for multi-step tasks
\end{itemize}

\section{Conclusion}

This project taught us how to implement and improve reinforcement learning algorithms. Starting from weak baselines, we systematically tested improvements and learned what works:

\textbf{PPO achieved 8.61\%} (69.5\% improvement) by using better hyperparameters, curiosity-driven exploration, and a larger network.

\textbf{DQN achieved [Partner fills]}.

Key lessons:
\begin{itemize}
    \item Sparse rewards need special exploration strategies
    \item Literature values are good starting points
    \item Documenting failures helps understand algorithms
    \item More training isn't always better
    \item Always evaluate separately from training
\end{itemize}

\textit{[Add final comparison of which approach worked better for Crafter and why]}

Code available at: \textit{[GitHub link]}

\begin{thebibliography}{9}

\bibitem{schulman2017}
J. Schulman et al., ``Proximal policy optimization algorithms,'' 2017.

\bibitem{mnih2015}
V. Mnih et al., ``Human-level control through deep reinforcement learning,'' Nature, 2015.

\bibitem{hafner2021}
D. Hafner, ``Benchmarking the spectrum of agent capabilities,'' 2021.

\bibitem{pathak2017}
D. Pathak et al., ``Curiosity-driven exploration by self-supervised prediction,'' 2017.

\bibitem{burda2019}
Y. Burda et al., ``Exploration by random network distillation,'' 2019.

\bibitem{schulman2016}
J. Schulman et al., ``High-dimensional continuous control using GAE,'' 2016.

\end{thebibliography}

\begin{figure}[t]
\centering
\includegraphics[width=0.48\textwidth]{fig2_achievement_comparison.pdf}
\caption{Achievement unlock rates across PPO evaluations. Shows improvement in tool crafting and discovery of rare achievements like coal and skeleton combat.}
\label{fig:achievements}
\end{figure}

\begin{figure}[t]
\centering
\includegraphics[width=0.48\textwidth]{fig3_summary_metrics.pdf}
\caption{Three key metrics (Score, Reward, Episode Length) across evaluations. All three improved with our changes.}
\label{fig:metrics}
\end{figure}

\begin{figure}[t]
\centering
\includegraphics[width=0.48\textwidth]{fig5_improvement_breakdown.pdf}
\caption{How much each evaluation contributed to the total 69.5\% improvement from baseline (5.08\%) to final result (8.61\%).}
\label{fig:breakdown}
\end{figure}

\end{document}
